\chapter{The compiler}\label{ch:compiler}

\begin{figure}[h]
\centering
\begin{tikzpicture}[node distance=5mm and 5mm, every node/.style={minimum height=8mm}]
  \node[box] (t) {term};
  \node[box, right=of t] (s) {sorts and\\shapes};
  \node[box, right=of s] (r) {rewrite\\($\varepsilon=0$)};
  \node[box, right=of r] (l) {lower\\cut sweep};
  \node[box, right=of l] (k) {KIR};
  \node[box, below=8mm of k] (sel) {selection\\per ISA};
  \node[box, left=of sel] (live) {liveness};
  \node[core, left=of live] (ra) {linear scan,\\no spill};
  \node[gate, left=of ra] (chk) {checker\\(from Lean)};
  \node[box, left=of chk] (enc) {encoders};
  \draw[edge] (t)--(s); \draw[edge] (s)--(r); \draw[edge] (r)--(l); \draw[edge] (l)--(k);
  \draw[edge] (k)--(sel); \draw[edge] (sel)--(live); \draw[edge] (live)--(ra); \draw[edge] (ra)--(chk); \draw[edge] (chk)--(enc);
  \draw[edge, dashed] (ra.north) -- (l.south);
\end{tikzpicture}
\caption{The compilation pipeline. An allocation that runs out of registers is never spilled: the
compiler goes back (dashed) and retries with a smaller group factor, fewer rows per iteration, or a cut of the fused region.}
\end{figure}

\section{A portable kernel IR}

Kernels (\texttt{Vapor.KIR.Kernels}: fused element-wise strips, reductions, matrix--vector products in
f32, bf16 and 4-bit formats, int8 matrix products, and the model operators) are written once over
virtual registers with \emph{types}: \texttt{:strip}, \texttt{:f16l} (sixteen f32 lanes),
\texttt{:i32acc} and others. Each backend sizes the types. The \textbf{group factor} $g$ generalises
the RVV register-group multiplier (LMUL) to every instruction set: an AVX2 strip with $g=4$ is four
\texttt{ymm} registers in lockstep, an RVV strip with $g=8$ is an \texttt{m8} group.

\begin{table}[h]
\centering\small
\begin{tabular}{@{}lllll@{}}
\toprule
type & RVV & AVX2 & AVX-512 & NEON \\
\midrule
\texttt{:strip} & m$g$ & $g$ ymm & $g$ zmm & $g$ q \\
\texttt{:f16l} & m4, vl$=16$ & 2 ymm & 1 zmm & 4 q \\
\texttt{:i32acc} & m$4g$ & $g$ ymm & $g$ zmm & 2 q \\
\bottomrule
\end{tabular}
\caption{One kernel IR, four register files.}
\end{table}

Selection happens \emph{before} allocation, as in production compilers. Lane-wise instructions are
emitted as bundles (one allocation unit), so a destination can reuse a source register that dies at
that point. On AVX-512, strip tails are one masked pass (\texttt{bzhi} into mask registers, masked loads
and stores only) instead of a scalar loop, and four-byte constants use embedded broadcast; the
reduction tree is lane for lane the same as AVX2's, so the bits are the same.

\section{Allocation without spill, checked}

\texttt{Vapor.KIR.Liveness} computes liveness by fixed point over the control-flow graph (values live on
a back edge cover the whole loop), with early-clobber constraints for the overlap rules of RVV's
widening instructions. \texttt{Vapor.KIR.RegAlloc} is a linear scan with aligned blocks of one, two,
four or eight registers.

There is no spill path. When registers run out, the allocator returns the pressure point, and the
compiler tries a smaller group factor, then fewer rows per iteration for the 4-bit products, then a
cut of the fused region at that point. Every accepted allocation is revalidated by
\texttt{check\_alloc/3}, which is extracted from Lean together with its soundness theorem
(\texttt{checkAlloc\_sound}): groups are aligned, inside the register file, outside the reserved
registers, and no two simultaneously live values share a register. Every emitted region is therefore
free of spill and clobber by construction, not by trust in the heuristic.

\section{Encoders}

\begin{table}[h]
\centering\small
\begin{tabularx}{\textwidth}{@{}lX@{}}
\toprule
target & encoding \\
\midrule
x86-64 (AVX2, AVX-512) & REX, three-byte VEX, EVEX (always \texttt{disp32}), ModR/M, SIB; System~V ABI, callee-saved registers saved only when used \\
AArch64 (NEON) & A64 words; AAPCS64, \texttt{x18} never allocated \\
RV64GCV & \texttt{vtype} fields, vector loads and stores with width fields, static round-to-nearest for scalar FP; full psABI; long branches as an inverted branch over a \texttt{jal} \\
SPIR-V (Vulkan) & a symbolic assembler that interns types and constants and orders the logical sections; \texttt{NoContraction} on every multiply and add under the canonical policy \\
MSL, StableHLO & emitted for Metal and for PJRT devices, and admitted only through the substrate airlock \\
\bottomrule
\end{tabularx}
\caption{Encoders. The product never calls an assembler.}
\end{table}

The tests validate every emitted instruction (every kernel constructor, in every backend, at every
group factor and policy) against GNU binutils, and every SPIR-V module against \texttt{spirv-val}.
Those tools are used only to check; the product encodes bytes itself.

\section{Cost model}

The arbiter predicts the time of a schedule on a substrate as
\[
  t = \sum_{\text{loops}} \max\!\left(\frac{F}{\pi}, \frac{Q}{\beta}, \frac{I}{\iota}\right) + \text{overheads},
\]
with floating-point work $F$ and bytes $Q$ counted from the schedule and $I$ the static count of
instructions in each hot loop as emitted, times its iterations. The third ceiling makes the model
honest: the 4-bit matrix--vector product is bound by instruction issue (about 1.1 instructions per
weight), not by bandwidth, and a two-ceiling model misses it by a factor of about five. Profiles are
declared per vendor; the dispatch decision is the argmin of the predicted time, taken on each node from
the certified work and the node's own profile.
